% SEGMENT OLP-0038-S01
% Part:sets-functions-relations
% Chapter: size-of-sets
% Section: enumerations-alt

\documentclass[../../../include/open-logic-section]{subfiles}

\begin{document}

\olfileid{sfr}{siz}{enm-alt}

\olsection{பட்டியலாக்கங்களும் \usetoken{S}{enumerable} கணங்களும்}

\begin{editorial}
  கணக்கோட்பாட்டில் செய்வதுபோல, இந்தப் பிரிவு பட்டியலாக்கங்களை
  $\Nat$ இன் (தொடக்கத் துண்டுகளுடனான) இருபுறச் சார்புகளாக
  வரையறுக்கிறது.
  எனவே \olref[enm]{sec} இன் வரையறைகளுடன் சிறிதளவு முரண்படுகிறது;
  அங்குள்ள எடுத்துக்காட்டுகள் அனைத்தையும் மீண்டும் கூறுகிறது.
  அந்தப் பிரிவைவிட இதன் விளக்கமும் சற்றுச் சுருக்கமானது.
\end{editorial}

% SEGMENT OLP-0038-S02
ஒரு முடிவுறு கணத்தின் !!{element}s[acc] பட்டியலிடுவதன் மூலம்
அதனை எளிதாகக் குறிப்பிடலாம். எடுத்துக்காட்டாக ஒரு கணத்தைப்
பின்வருமாறு வரையறுக்கும்போது இதையே செய்கிறோம்:
\[
  A = \{a_1, a_2, \ldots, a_n\}.
\]
!!{element}s $a_1$, \dots, $a_n$ அனைத்தும் வெவ்வேறானவை எனக் கொண்டால்,
இது $A$ க்கும் முதல் $n$ இயல் எண்களான
$0$, \dots, $n-1$ க்கும் இடையில் !!a{bijection} ஐத் தருகிறது.
மறுதிசையில், ஒவ்வொரு முடிவுறு கணத்திலும் முடிவுறு எண்ணிக்கையிலான
!!{element}s மட்டுமே இருப்பதால், ஒவ்வொரு முடிவுறு கணத்தையும்
இத்தகைய ஒத்திசைவில் வைக்கலாம்.
வேறு சொற்களில், $A$ முடிவுறு கணமாக இருந்தால்,
$A$ க்கும் $\{0, \dots, n-1\}$ க்கும் இடையில் !!a{bijection} உள்ளது;
இங்கு $n$ என்பது $A$ இன் !!{element}s எண்ணிக்கை.

சில வகை முடிவுறாத கணங்களையும் அனுமதித்தால்,
அவற்றில் சிலவற்றையும் பட்டியலிட அனுமதிப்போம்.
ஒரு முடிவுறாத கணத்திற்கும் முழு $\Nat$ க்கும் இடையில்
!!a{bijection} இருப்பதையே அந்தக் கணத்தின் பட்டியலாக்கம் என்று கூறி,
இதனைத் துல்லியமாக்கலாம்.

\begin{defn}[பட்டியலாக்கம் — கணக்கோட்பாட்டு வரையறை]
ஒரு கணம் $A$ இன் \emph{பட்டியலாக்கம்} என்பது,
அதன் வீச்சகம் $A$ ஆகவும் அதன் சார்பகம்
$\{0, 1, \ldots, n\}$ என்ற இயல் எண்களின் தொடக்கக் கணமாகவோ
முழு இயல் எண்களின் கணம் $\Nat$ ஆகவோ உள்ள !!a{bijection} ஆகும்.
\end{defn}

% SEGMENT OLP-0038-S03
\begin{explain}
\emph{பட்டியலாக்கம்} என்ற சொல்லின் இந்தப் பயன்பாட்டிற்கு
ஓர் உள்ளுணர்வு அடிப்படை உள்ளது.
ஒரு கணம் $A$ ஐப் பட்டியலிட்டுள்ளோம் என்று சொல்வது,
$A$ இன் உறுப்புகளை எண்ணிச் செல்ல உதவும் !!a{bijection} $f$
உள்ளது என்று சொல்வதாகும்.
$0$ ஆவது உறுப்பு $f(0)$, $1$ ஆவது $f(1)$, \ldots,
$n$ ஆவது $f(n)$, \ldots{} ஆகும்.
\footnote{ஆம், $0$ இலிருந்து எண்ணுகிறோம்.
நிச்சயமாக ஒன்றிலிருந்தும் தொடங்கலாம். அதனால் பெரிய வேறுபாடு ஏதுமில்லை;
$\Nat$ க்குப் பதிலாக $\PosInt$ ஐ மட்டும் பயன்படுத்த வேண்டியிருக்கும்.}
பின்வருவதைச் சேர்ப்பதன் மூலம் இதற்கான காரணத்தை மேலும் தெளிவாக்கலாம்:
\end{explain}

% SEGMENT OLP-0038-S04
\begin{defn}
  \ollabel{defn:enumerable}
  ஒரு கணம் $A$ என்பது $A = \emptyset$ ஆக இருக்கும்போது அல்லது
  $A$ க்கு ஒரு பட்டியலாக்கம் இருக்கும்போது, அப்போதும் அப்போது மட்டுமே,
  !!{enumerable} ஆகும்.
  $A$ ஆனது !!{enumerable} ஆக இல்லாதபோது, அப்போதும் அப்போது மட்டுமே,
  $A$ என்பது !!{nonenumerable} என்கிறோம்.
\end{defn}

% SEGMENT OLP-0038-S05
\begin{explain}
எனவே ஒரு கணம் வெற்றுக் கணமாக இருக்கும்போது அல்லது
ஒரு பட்டியலாக்கத்தைக் கொண்டு அதன் !!{element}s[acc]
எண்ணிச் செல்ல முடியும் போது, அப்போதும் அப்போது மட்டுமே,
அது !!{enumerable} ஆகும்.
\end{explain}

% SEGMENT OLP-0038-S06
\begin{ex}
இயல் எண்களைப் பட்டியலாக்கும் ஒரு சார்பு,
$\Id{\Nat}(n) = n$ எனக் கொடுக்கப்படும்
ஒன்றாதல் சார்பு $\Id{\Nat} \colon \Nat \to \Nat$ ஆகும்.
\emph{நேர்ம} இயல் எண்களான $\Nat^+ =
\Nat \setminus \{0\}$ ஐப் பட்டியலாக்கும் சார்பு,
$g(n) = n + 1$ என்ற தொடரிச் சார்பாகும்.
\end{ex}

% SEGMENT OLP-0038-S07
\begin{prob}
ஒரு கணம் $A$ என்பது $A = \emptyset$ ஆக இருக்கும்போது அல்லது
!!a{surjection} $f\colon \Nat \to A$ இருக்கும்போது,
அப்போதும் அப்போது மட்டுமே, !!{enumerable} ஆகும் எனக் காட்டுக.
!!a{injection} $g\colon A \to \Nat$ இருக்கும்போது,
அப்போதும் அப்போது மட்டுமே, $A$ என்பது !!{enumerable} ஆகும் எனக் காட்டுக.
\end{prob}

% SEGMENT OLP-0038-S08
\begin{ex}
பின்வருமாறு கொடுக்கப்படும்
$f\colon \Nat \to \Nat$, $g \colon \Nat \to \Nat$ ஆகிய சார்புகள்:
\begin{align*}
  f(n) & = 2n \text{ மற்றும்}\\
  g(n) & = 2n+1
\end{align*}
முறையே இரட்டை இயல் எண்களையும் ஒற்றை இயல் எண்களையும்
பட்டியலாக்குகின்றன.
ஆனால் இவற்றில் எதுவும் !!{surjective} பண்புடையதல்ல;
எனவே எதுவும் $\Nat$ இன் பட்டியலாக்கமல்ல.
\end{ex}

% SEGMENT OLP-0038-S09
\begin{prob}
வர்க்க எண்கள் $1$, $4$, $9$, $16$, \dots ஆகியவற்றின்
ஒரு பட்டியலாக்கத்தை வரையறுக்கவும்.
\end{prob}

% SEGMENT OLP-0038-S10
\begin{ex}
$\lceil x \rceil$ என்பது $x$ ஐ மிக அருகிலுள்ள முழுவுக்கு
மேல்நோக்கி முழுமைப்படுத்தும் \emph{மேல்முழுச்} சார்பு என்க.
பின்வருமாறு கொடுக்கப்படும் $f \colon \Nat \to \Int$ என்ற சார்பு:
\[
  f(n) = (-1)^{n} \left\lceil\tfrac{n}{2}\right\rceil
\]
முழுக்களின் கணம் $\Int$ ஐப் பின்வருமாறு பட்டியலாக்குகிறது:
\[
\begin{array}{c c c c c c c c}
f(0) & f(1) & f(2) & f(3) & f(4) & f(5) & f(6) & \dots \\ \\
\big\lceil \tfrac{0}{2} \big\rceil & -\big\lceil \tfrac{1}{2}\big\rceil &  \big\lceil \tfrac{2}{2} \big\rceil & -\big\lceil \tfrac{3}{2} \big\rceil & \big\lceil \tfrac{4}{2} \big\rceil  & -\big\lceil \tfrac{5}{2}\big\rceil & \big\lceil \tfrac{6}{2} \big\rceil & \dots \\ \\
0 & -1 & 1 & -2 & 2 & -3 & 3& \dots
\end{array}
\]
நேர்ம, எதிர்ம முழுக்களுக்கு இடையில் முன்னும் பின்னும் ``தாவி'',
$\Int$ இன் மதிப்புகளை $f$ எவ்வாறு உருவாக்குகிறது என்பதைக் கவனிக்கவும்.
$f$ பின்வருமாறு வகைப்பிரிவுகளால் வரையறுக்கப்பட்டதாகவும் கருதலாம்:
\[
f(n) = \begin{cases}
  \frac{n}{2} & \text{$n$ இரட்டை எனில்}\\
  -\frac{n+1}{2} & \text{$n$ ஒற்றை எனில்}
  \end{cases}
\]
\end{ex}

% SEGMENT OLP-0038-S11
\begin{prob}
$A$, $B$ ஆகியவை !!{enumerable} எனில்,
$A \cup B$ என்பதும் எண்ணிடத்தக்கது எனக் காட்டுக.
\end{prob}

% SEGMENT OLP-0038-S12
\begin{prob}
$A_1$, $A_2$, \dots, $A_n$ ஆகிய அனைத்தும் !!{enumerable} எனில்,
$A_1 \cup \dots \cup A_n$ என்பதும் எண்ணிடத்தக்கது என்பதை,
$n$ மீதான தொகுத்தறிதல் மூலம் காட்டுக.
\end{prob}

\end{document}
